% agents-life-cycle.tex — Chapter 8 of the thesis: the life cycle of the concept-agents.
%
% Sources: CONCEPT_AGENT_LIFECYCLE.md (the normative contract), MODEL_HISTORY.md (stage 4 theory),
% and LIFECYCLE.tex (the first, shorter draft, which this supersedes).
% Reads the contract against the factorization of Chapter 7 (FACTORIZATION_ANCHORS.tex).
%
% Build standalone:  pdflatex agents-life-cycle.tex   (twice, for cross-references)
% To include in main.tex: strip the preamble down to \chapter and drop the \setcounter.
%
\documentclass[11pt,a4paper]{report}
\usepackage[T1]{fontenc}
\usepackage{lmodern}
\usepackage[margin=2.6cm]{geometry}
\usepackage{amsmath,amssymb,amsthm}
\usepackage{empheq}
\usepackage{pifont}
\usepackage{array}
\usepackage{booktabs}
\usepackage{enumitem}
\usepackage[hidelinks]{hyperref}
\usepackage{microtype}

\newcommand{\given}{\,|\,}
\newcommand{\R}{\mathbb{R}}
\newcommand{\stage}[1]{\textsc{#1}}
\newcommand{\Create}{\stage{create}}
\newcommand{\Update}{\stage{update}}
\newcommand{\Remove}{\stage{remove}}
\newcommand{\History}{\stage{history}}
\newcommand{\Ownership}{\stage{ownership}}

\begin{document}

\setcounter{chapter}{7}
\chapter{The life cycle of the concept-agents}

\begin{flushleft}
\textit{No\'e Zapata}
\end{flushleft}

\vspace{4pt}

\section{Introduction}
\label{sec:lc-intro}

The previous chapter wrote down the distribution the agents approximate. It says what an instance
\emph{is} --- a Gaussian over six geometric parameters and a log-odds of existence --- and it says
how the measurements bear on it. It does not say when an instance should come into being, when it
should stop existing, or what a process is allowed to do to the shared graph in the meantime. Those
questions are answered by a second document, written independently and for a different reason.

Every \texttt{<obj>\_concept} agent must satisfy a normative behavioural contract, stated once in
\texttt{CONCEPT\_\allowbreak AGENT\_\allowbreak LIFECYCLE.md} and audited per agent. Its five stages
are \Create, \Update, \Remove, \History{} and \Ownership. It is not a description of what the agents
do; it is a specification of what they are permitted to do, and several agents are in breach of it
at the time of writing (\S\ref{sec:audit}).

That contract was written from experience rather than from theory. It exists because the same defect
class was found and fixed three separate times in three different agents --- an existence channel
voting to delete objects through a wall, in \textsf{refrigerator}, then \textsf{door}, then
\textsf{table} and \textsf{cabinet} --- and because a threshold that removed objects the sensor
provably could not see survived inside a hand-rolled copy of a stage that already had a shared,
correct implementation. Both were possible because the life cycle was structure without a contract.
Its purpose is that a stage stated once is a stage a new agent cannot forget.

The two documents therefore approach one system from opposite directions: Chapter~7 from the joint
distribution and the seven approximations that make it computable, this chapter from the failures
and the rules written to prevent them. They were produced independently, which is what makes their
comparison worth something. Where they agree, the contract acquires a derivation and stops being a
list of house rules. Where they disagree, one of them is missing a piece --- and in this case both
are, in different places.

The result of the comparison is the structure of this chapter. Four of the five stages restate
factors or assumptions of Chapter~7, and one of them, \History, improves on the assumption as
stated (\S\ref{sec:mapping}). The fifth, \Create, has no counterpart in the model at all: birth adds
a term to a product that contains nothing deciding it should be added (\S\ref{sec:create}). \Update{}
carries a hard admissibility gate that looks like a second gap of the same kind and, on closer
reading, is not --- the one exactly-solvable factor of Chapter~7 rules against the continuous
alternative, and the gate turns out to be a runtime defence of assumption A2
(\S\ref{sec:update}). \Remove{} contributes a correction \emph{back} to Chapter~7: the defect found
three times is a missing argument in one factor of the joint (\S\ref{sec:remove}). \History{} is a
single move --- a constant becomes an inferred latent --- applied to the two constants that govern
accumulation (\S\ref{sec:history}). \Ownership{} is an assumption of the factorization wearing
administrative clothes, and its failures are the failures of an assumption that no process is
enforcing (\S\ref{sec:ownership}).

\section{What Chapter 7 leaves on the table}
\label{sec:recap}

Three results of the previous chapter are used throughout, and are restated here so this chapter can
be read on its own.

\paragraph{The seven assumptions.} With $\theta_i\in\R^6$ the geometry of instance $i$, $e_i$ its
existence log-odds, $\mathcal{I}_c$ the instances of class $c$, $M_t$ the per-cycle mask-to-instance
matching, $X$ the trajectory and $R$ the room polygon:

\begin{center}
\begin{tabular}{@{}l l p{8.6cm}@{}}
\toprule
& Acts on & What it assumes \\
\midrule
A1 & localisation & objects and pose are updated alternately, each holding the other fixed \\
A2 & masks & the trajectory marginal factorises frame by frame \\
A3 & matching & one matching $\hat M$, not a distribution over matchings \\
A4 & masks & object classes are independent --- one agent each \\
A5 & masks & instances within a class are independent \\
A6 & masks & geometry and existence factorise: $q(\theta_i)\,q(e_i)$ \\
A7 & masks & the state is Markov in time \\
\bottomrule
\end{tabular}
\end{center}

\paragraph{The free energy an agent minimises.} Per instance,
\begin{equation}
F_i\;=\;
\underbrace{\mathbb{E}_{q}\!\left[-\ln P\!\left(\bar Z_i\given\theta_i,\hat X\right)\right]}_{\text{accuracy}}
\;+\;
\underbrace{\mathrm{KL}\!\left[\,q(\theta_i)\,\|\,P(\theta_i\given e_i)\,\right]}_{\text{complexity}} ,
\label{eq:lc-F}
\end{equation}
where the accuracy term is the negative log of a \emph{mixture}: each mask point is explained by the
object's surface, by a uniform clutter component, or by a known wall. The clutter component's weight
is a configured constant; \S\ref{sec:pfa} argues it should not be.

\paragraph{The ceiling.} The single factor of Chapter~7 evaluated without approximation is the
per-frame marginalisation of the localisation error $\delta_t$. Stacking $N$ point residuals of one
mask with Jacobian $G$ with respect to that shared offset,
\begin{equation}
\mathcal{I}_t \;=\; G^{\!\top}\!\left(\sigma^2 I+G\Sigma_\delta^t G^{\!\top}\right)^{-1}\!G
\;=\;
\left(\sigma^2\left(G^{\!\top}G\right)^{-1}+\Sigma_\delta^t\right)^{-1}
\;\xrightarrow[\ N\to\infty\ ]{}\;
\Lambda_c \;:=\; \left(\Sigma_\delta^t\right)^{-1}.
\label{eq:lc-sat}
\end{equation}
One frame's information about $\theta_i$ saturates at $\Lambda_c$, however many points it carries.
Two clauses of the contract are consequences of \eqref{eq:lc-sat} --- one in \Update{}
(\S\ref{sec:whygate}) and one in \Remove{} (\S\ref{sec:saturation}) --- and neither was derived that
way. They were found by debugging, twice, in two different channels.

\section{Where the five stages land in the model}
\label{sec:mapping}

\begin{center}
\begin{tabular}{@{}l l p{8.0cm}@{}}
\toprule
Stage & In the model & \\
\midrule
\Create & \textbf{nothing} &
  birth adds an index to $\mathcal{I}_c$, and therefore a term to $\prod_{i\in\mathcal{I}_c}$.
  No factor decides that the term should be added. \S\ref{sec:create} \\[4pt]
\Update & accuracy term; A2, A5, A7 &
  ``a frame must earn the right to move geometry'' is the admissibility side of the measurement
  likelihood; the module the contract mandates is the recursive-Laplace engine itself. The gate
  itself defends A2. \S\ref{sec:update} \\[4pt]
\Remove & $q(e_i)$, A6; visibility &
  ``evidence, never a miss-counter'' requires removal to be driven by the existence factor, which
  A6 is what makes separable from geometry in the first place. Its evidence comes from the
  visibility factor Chapter~7 declares inert. \S\ref{sec:remove} \\[4pt]
\History & A7; the clutter rate &
  ``no constant may stand in for experience'' refines the Markov transition --- it forbids a constant
  process noise $Q$ --- and applies the same move to the mixture's false-alarm weight.
  \S\ref{sec:history} \\[4pt]
\Ownership & A4 &
  one agent owns the instances of one class, which is well defined only because A4 made the classes
  independent. \S\ref{sec:ownership} \\
\bottomrule
\end{tabular}
\end{center}

\noindent
\History{} is the clearest case of the contract being ahead of a naive reading of the model. A7 as
stated permits any transition $P(\theta_i^t\given\theta_i^{t-1})$, and the deployed choice --- a
constant process noise applied every frame, observed or not --- satisfies it. The contract observes
that such a $Q$ inflates covariance on the agent's clock rather than on evidence, and requires the
growth to depend on what actually happened. That is a strictly better model, arrived at from a
symptom rather than from the algebra.

\Ownership{} is the opposite case: an assumption of the factorization that has no probabilistic
content left in it by the time it reaches the code, only a configuration file that can be wrong.

\section{\Create: the stage with no factor}
\label{sec:create}

\subsection{What the question actually is}

Deciding whether an object exists is a comparison between two models of \emph{different order}: the
world with $n$ instances of this class and the world with $n+1$. Under the variational treatment of
Chapter~7 that comparison has a definite form. Writing $F^{(n)}$ for the total free energy of the
class under $n$ instances,
\begin{align}
\Delta F \;=\; F^{(n+1)}-F^{(n)}
\;=\;&
\underbrace{\mathbb{E}_{q}\big[-\ln P(\bar Z\given\theta_{1:n+1})\big]
          -\mathbb{E}_{q}\big[-\ln P(\bar Z\given\theta_{1:n})\big]}_{\text{accuracy gained by explaining those points}}
\nonumber\\[4pt]
&+\;
\underbrace{\mathrm{KL}\!\left[\,q(\theta_{n+1})\,\|\,P(\theta_{n+1}\given e)\,\right]}_{\text{complexity paid for the new object}} .
\label{eq:lc-deltaF}
\end{align}
The new instance is warranted exactly when $\Delta F<0$: when the measurements it claims are
explained better by a new object than by the clutter component and by the instances already present,
by more than the cost of instantiating a new set of parameters. The second term is not a penalty
bolted on for regularisation --- it is the same complexity term that already appears in
\eqref{eq:lc-F} for every live instance, and it is what stops the model from explaining every stray
mask with a new refrigerator.

Nothing in the assembled posterior of Chapter~7 computes \eqref{eq:lc-deltaF}, because
$\mathcal{I}_c$ enters that posterior as a \emph{fixed} index set. This is a choice, not a necessity:
the random-finite-set formulations of multi-target tracking carry a birth intensity as an explicit
term of the model \cite{mahler}, and the classical alternative --- M-of-N track initiation
\cite{barshalom} --- is a surrogate of exactly the kind described next, adopted for exactly the same
reason.

\subsection{What is computed instead}

\paragraph{Birth is tracker-only.} The contract's first invariant is architectural: no agent may
birth an instance from its own logic. Birth happens in one shared module,
\texttt{common/instance\_tracker}, and an agent must never create a graph node for an instance the
tracker has not birthed. This is what stops the birth criterion from being re-derived, differently,
eight times.

\paragraph{Birth is evidence-accumulating, not a counter.} A detection matures a \emph{pending
candidate} over \texttt{birth\_frames}, and each frame contributes a weight,
\texttt{Detection::birth\_evidence}, defaulting to $1.0$. That weight is the only sanctioned hook for
a birth prior, and it is continuous: a corroborated detection matures faster, a suspect one slower.
Two further rules make the stage well-posed in a multi-sensor system: a detection may be marked
\emph{non-birthable}, allowing it to refine an existing track but never to start one --- one sensor
births, peripheral sensors only cue attention --- and a candidate must stand clear of
\texttt{birth\_min\_sep\_m} from every existing track \emph{and} every other pending candidate.

In the refrigerator agent the per-frame weight is a product of several factors, each a proxy for a
term of \eqref{eq:lc-deltaF}:

\begin{center}
\begin{tabular}{@{}l p{9.4cm}@{}}
\toprule
Factor & What it stands in for \\
\midrule
detector confidence & the likelihood's weight on the class label \\
range attenuation   & measurement precision, which falls with distance \\
frame admissibility & whether this frame may move geometry at all (\S\ref{sec:update}) \\
shape plausibility  & the prior term $P(\theta_{n+1}\given e)$, evaluated on the raw cloud \\
room-envelope support & the same prior, evaluated against the wall polygon \\
residual surprise (optional) & unexplained occupancy under the detection: the accuracy term \\
\bottomrule
\end{tabular}
\end{center}

\noindent
The list is a reasonable decomposition. Every entry corresponds to a term that would appear in
$\Delta F$, and the shared implementation means all agents birth on the same terms.

\subsection{What the surrogate costs}

Three things are lost by accumulating a weighted count rather than evaluating
\eqref{eq:lc-deltaF}.

\paragraph{The units are gone.} A streak is a dimensionless number compared against a frame count.
$\Delta F$ is in nats, and can therefore be traded against any other quantity in the model --- the
existence log-odds, a merge decision, the expected free energy of a viewpoint. Because the birth
score is not in nats, birth cannot participate in any of those trades, and the boundary it is
compared against is a threshold with no principled value.

\paragraph{Birth and death are not the same currency.} Removal is driven by an existence log-odds
accumulated from occupancy carve and silhouette evidence (\S\ref{sec:remove}). Birth is driven by a
streak. The two decisions concern the same proposition --- does this object exist --- and are taken
in incommensurable units, so an object can be born on evidence that would not have kept it alive,
and the asymmetry is invisible to both mechanisms. The visible symptom is a whole corrective
subsystem: the refrigerator agent carries shape-plausibility evidence routed into the existence
channel plus a soft singleton inhibition, so that a strongly-supported instance suppresses weaker
duplicates. It works, and it is carefully built. It is also a loop whose sole purpose is to undo
births that should not have happened, and its own criterion is, again, in different units from the
one that created the object.

\paragraph{The comparison is available and not made.} $\Delta F$ is not out of reach. A candidate
that has matured has an associated set of measurements; instantiating a provisional instance, fitting
it, and reading off the resulting free energy against the null is the same computation the agent
already performs every cycle for its live instances. What is missing is the plumbing, not the theory.
A birth gate that computed \eqref{eq:lc-deltaF} on the candidate's accumulated measurements would
subsume the plausibility multiplier, the envelope support, the residual-surprise fusion and the
retraction subsystem, because all four are estimating parts of the same scalar.

\subsection{The model has no birth time either}
\label{sec:provenance}

There is a second, quieter absence at this stage. The posterior of Chapter~7 is a distribution over
what exists; it has no notion of \emph{when} an index entered the product, and the shared graph
inherits that silence. The CRDT's own per-attribute timestamps do not fill it: they move on every
write, so they date the last update, not the birth.

The contract therefore requires the missing quantity to be recorded outside the model. Every node
this fleet inserts carries its birth time, written once, immediately before insertion, and never
refreshed --- as a millisecond epoch for arithmetic and as a local civil ISO-8601 string for the
graph viewer. The scope is every node, not just instances: affordance children, room, wall and floor
nodes, the controller's temporary obstacles, the retina's channel nodes. A stale node in the
shared graph can then always be dated against the run that leaked it.

One rule follows from taking the model seriously: a re-acquired object gets a \emph{new} stamp. It is
a new node, with a new identifier, and only its name carried over. ``When did this object first
appear, across deaths'' is a question about $e_i$ over time --- an existence-belief question --- not
a question about node provenance, and conflating the two would put a persistent identity into a
factor that does not have one.

\subsection{Why the gap matters more in this variant}
\label{sec:variant}

Chapter~7 is the variant in which converged objects are published back to the localiser as SE(2)
landmarks. That closes a loop the baseline cuts, and it has a consequence for this stage that is
worth stating plainly: with the loop closed, the cost of a bad birth stops being local. A phantom
instance perturbs $\hat X$, and $\hat X$ enters every other agent's $\Sigma_\delta^t$ through
\eqref{eq:lc-sat}. A5 --- instances are independent because they own disjoint measurement sets ---
is no longer even approximately true, since all of them now share a pose.

The birth criterion is the fleet's only defence against that, and it is the one stage with no
principled criterion. A system that feeds objects back into localisation should be \emph{more}
demanding about model-order selection than one that does not, and this one is currently less. That
asymmetry is the strongest argument in this chapter for closing the \Create{} gap first, and it is
also part of why the room-side landmark switch is off today.

\section{\Update: a frame must earn the right to move geometry}
\label{sec:update}

The required module is the recursive-Laplace engine, \texttt{rc::ai::update}, driven by the agent's
own generative model. The Bayesian mathematics lives in the engine exactly once; an agent supplies a
residual and its Jacobian, and nothing else. That much is a straight restatement of
\eqref{eq:lc-F}.

The contract then adds something the model does not contain. A frame may touch the geometric mean
only when it is

\begin{enumerate}[label=\arabic*.,leftmargin=1.9em,itemsep=3pt]
\item \textbf{resolvable} --- enough un-cluttered surface to constrain the state, judged on point
      mass, clutter fraction and truncation. \emph{Not} distance: a dense far view beats a starved
      near one;
\item \textbf{centred} --- the object lies in the sensor's reliable region, not clipping the
      periphery;
\item \textbf{still} --- ego-motion smear is below the level at which the mask is a shared blur.
\end{enumerate}

\noindent
Otherwise the cycle is \texttt{predict()}-only: the mean is held and $\Sigma_{\theta_i}$ inflates.
This is a hard predicate --- a boundary, in a system whose stated preference is for continuous
precisions --- and it is flagged as such in the contract.

\subsection{Why a covariance cannot do this job}
\label{sec:whygate}

The obvious objection is that a poor frame should carry a large observation covariance and lose to
the prior on its own, continuously and without a boundary. This is the repair the project applies
everywhere else, and it is the one an earlier draft of this chapter recommended. It does not work
here, and \eqref{eq:lc-sat} is why.

Suppose an inadmissible frame carries not merely extra noise but a systematic offset $b$: motion
smear displaces the mask's centroid in the direction of travel, truncation biases the fitted extent
towards the visible side, an off-centre object is deprojected with the periphery's distortion. Let
the prior information be $\lambda_0$, and let $n$ such frames each contribute information $\lambda$
about the affected degree of freedom. The posterior mean is displaced by
\begin{equation}
\Delta\mu_n \;=\; \frac{n\lambda}{\lambda_0+n\lambda}\;b
\;\xrightarrow[\ n\to\infty\ ]{}\; b .
\label{eq:lc-bias}
\end{equation}
Attenuation multiplies $\lambda$ by some $\alpha\in(0,1]$. It changes the \emph{rate} at which
\eqref{eq:lc-bias} approaches $b$ and does not change the limit, for any $\alpha>0$. Only $\alpha=0$
--- a gate --- bounds the displacement. And $\alpha$ cannot be driven arbitrarily close to zero by
the covariance model, because \eqref{eq:lc-sat} guarantees the frame's information saturates at
$\Lambda_c$, a \emph{nonzero} asymptote: the engine will not let a frame contribute more than
$\Lambda_c$, but neither does it let a downweighted frame contribute nothing. A bad frame is
attenuated and still moves the mean, and accumulation over hundreds of frames beats attenuation.

The argument turns entirely on $b$ being \emph{shared across the run of frames}. If the per-frame
offsets were independent and zero-mean, averaging would dispose of them at rate $n^{-1/2}$ and a
covariance would be the right and sufficient answer. They are not independent: the robot moves
smoothly, so consecutive frames smear the same way; an object off the optical axis stays off it for
a run of frames; a truncated view is truncated on the same side until the robot turns.

That is precisely the assumption A2 makes and the justification A2 offers. A2 marginalises the
localisation error frame by frame on the grounds that each frame is deprojected with its own
capture-stamp transform, so its offset is drawn afresh --- and Chapter~7 records that this is
\textbf{false across short windows}, measured: pooling eight frames' raw points into one estimate
inflated a refrigerator's footprint from $0.669$ to $0.915$\,m while \emph{lowering} its reported
$\sigma_H$ from $0.300$ to $0.190$. The inadmissible frames of this section are the same failure
arriving through a different door. They do not break the likelihood; they break the
independence A2 assumes between one frame's offset and the next.

\begin{quote}
The admissibility gate is not a missing term in the measurement model. It is a runtime guard on the
justification of A2, applied to the frames where that justification is known to fail.
\end{quote}

\noindent
This reading also says what the principled version would be, and it is not ``fold the three
covariates into $R$''. It is to give the correlated part of the offset its own nuisance variable ---
a per-window pose with a prior, marginalised as $\delta_t$ is marginalised --- and to make it
identifiable. Chapter~7 sets out the three requirements for that construction and concludes it is
not worth the cost for a six-parameter box with tight priors. Given that conclusion, the gate is the
correct engineering answer and not a compromise: it costs a few frames of intake to avoid a bias the
model has no coordinate for.

\subsection{What the gate does not touch}

The gate is deliberately narrow. Association and existence do \emph{not} read it: an inadmissible
frame still confirms an instance, still updates $\hat M$, and still feeds the existence channel.
Only the geometry moves are frozen.

In the terms of \eqref{eq:lc-F}, the gate suppresses the accuracy term of the geometry block for one
cycle and leaves the rest of the posterior running. This is not a detail of implementation. A gate
that also froze existence would turn every occluded or blurred stretch into evidence of absence,
which is precisely the failure \Remove{} spends its whole specification avoiding
(\S\ref{sec:remove}); and a gate that froze association would drop tracks across exactly the
manoeuvres where the robot is moving.

\subsection{Per-frame quality is a covariance, never a gate}

The narrowness cuts the other way too. Everything that is genuinely a \emph{precision} question ---
range, obliquity, coverage, ego-motion magnitude --- enters through the per-frame common-mode
$\Sigma_\delta^t$, Woodbury-marginalised in the engine as in \eqref{eq:lc-sat}, so that $N$
correlated points cannot collapse $\sigma$. New quality signals go there, not into a new \texttt{if}.
The gate is reserved for the one situation the covariance provably cannot express: a shared bias
across a run of frames.

\subsection{Observability covariates must not be circular}
\label{sec:circular}

One further rule belongs to this stage, and it is the subtlest in the contract.

A covariate that decides whether a degree of freedom is observable must not be computed from the
point estimate of that same degree of freedom. If it is, a wrong estimate declares its own error
unobservable and protects itself: the belief stops updating precisely because it is wrong, and the
mechanism that would have corrected it reads its own mistake as a reason not to look.

The live precedent is the chair agent. Its backrest-normal obliquity --- the angle at which the
sensor sees the backrest, and hence how much a frame can say about yaw --- was computed from the
belief's own yaw. With that yaw about $45^\circ$ wrong, the covariate reported
$\mathrm{oblq\_cos}=0.001$: from where the belief thought the backrest was, the backrest was almost
edge-on and the frame was declared uninformative. The yaw froze at the wrong value and stayed there.

Stated variationally, the rule is that such a covariate must be evaluated \emph{under the posterior}
$q(\theta_i)$, or from the data directly, rather than at $\mu_i$. A quantity that is a function of
what we do not yet know is a quantity whose uncertainty must be carried, and the cheap substitution
$f(\mathbb{E}_q[\theta])$ for $\mathbb{E}_q[f(\theta)]$ is exactly what removes the uncertainty that
would have kept the door open. The same structure appears again in \S\ref{sec:attribution}, in a
completely different stage, and the contract forbids it in both places for the same reason.

\subsection{The gate and the planner are one mechanism}

A hard admissibility predicate is only defensible in a system that can go and produce admissible
frames. On its own it is starvation: an object seen badly is never updated, and nothing changes
that.

The other half is already present. Chapter~7's expected-free-energy term is implemented as the
next-best-view planner, which scores each face of the fitted box by the D-optimal expected entropy
reduction on $\Sigma_{\theta_i}$ and proposes the standoff pose that observes the winner. The
covariance the gate inflates on a frozen cycle is the same covariance the planner reads to decide
where to look, so a run of inadmissible frames does not merely suppress intake --- it raises the
epistemic value of going to look properly, and the planner acts on it.

That is the sense in which the gate is the attention analogue the contract claims: intake is
suppressed \emph{outside} a fixation and sought \emph{inside} one, rather than being uniformly
down-weighted everywhere. Perception and action are scored in one currency, and the gate is what
makes the currency worth having.

\section{\Remove: evidence, never a miss-counter}
\label{sec:remove}

The required module is the shared existence belief, \texttt{common/existence\_belief}. One scalar
log-odds $L$ per instance --- this is $e_i$, and A6 is what makes it separable from the geometry ---
clamped to $\pm L_{\max}$ so that it stays finite and, more importantly, recoverable.

\subsection{One scalar, three verdicts}

\begin{center}
\begin{tabular}{@{}p{6.4cm} l p{5.1cm}@{}}
\toprule
Observation & Verdict & \\
\midrule
occupancy: a return, or a lit mask inside the volume & \textbf{confirms} & holds $L$ up \\[2pt]
predicted-visible but absent & \textbf{removes} & the only removal evidence there is \\[2pt]
occluded, out of field of view, or not probed & \textbf{holds} & never absence \\
\bottomrule
\end{tabular}
\end{center}

\noindent
The whole stage is contained in the third row. A miss is not evidence of absence; a miss
\emph{where the object would have been seen} is. Everything else in this section is the consequence
of taking that seriously.

\subsection{The missing argument}
\label{sec:walls}

``Where the object would have been seen'' is a visibility computation, and Chapter~7 has a factor for
it: $P(M\given\theta,e,X)$, the block that says what is visible from where. That block is declared
inert in the free energy, and correctly so --- A3 collapses $M$ to a delta, and a delta contributes
neither entropy nor expectation, so the matching is free in the literal sense that $F$ cannot see it.

But \Remove{} does not read $M$. It reads the \emph{visibility} that block encodes, and it reads it
as the weight on every removal increment. The factor Chapter~7 declares free is the one this stage
depends on entirely, and nobody was checking it, because in the free energy there was nothing to
check.

Now write the factor out and look at its arguments. It is a function of $\theta$, $e$ and $X$. It is
not a function of $R$ --- and the room polygon is exactly what tells you that the object is behind a
wall.

\begin{quote}
A wall is not a detector class. With visibility computed from masked occluders alone, an object in
the next room projects into the frustum, lands on unmasked wall pixels, finds no mask there, and
votes for its own removal \emph{at full strength}.
\end{quote}

\noindent
That is the defect found three times --- refrigerator, then door, then table and cabinet --- and it
is legible in the joint as a missing argument. The correct factor is
\begin{equation}
P\!\left(M\given\theta,e,X,\mathbf{R}\right),
\label{eq:lc-visR}
\end{equation}
and the shared implementation the contract mandates, \texttt{rc::occlusion::walls\_block()}, is
\eqref{eq:lc-visR} rather than its predecessor. This is the one place where this chapter amends the
previous one rather than deriving from it. It is also a small vindication of writing the joint down:
a factor whose argument list is wrong is a bug that recurs in every process that recomputes it, and
there is no amount of per-agent testing that finds it, because each agent's version is locally
self-consistent.

\subsection{Weighting on the wrong side of a ceiling}
\label{sec:saturation}

Absence must be weighted by $P(\text{detect}\given\text{present})$: a miss argues for removal only to
the degree the sensor could have resolved a present object from here. The clause is obvious. Its
implementation was wrong for a reason that is not.

The weight must scale the \emph{saturated} per-cycle log-odds increment, not the raw pixel or point
count that produced it. Scaling the count is a no-op. The per-cycle evidence hits its ceiling
$\pm\ell_{\mathrm{occ}}$ long before the count runs out, so even a $2\%$-resolvable view still lands
at the full increment, and the weight --- carefully computed, correctly motivated --- had no effect
whatsoever on the outcome.

This is \eqref{eq:lc-sat} again, in a channel that has nothing to do with geometry. The lesson
generalises past both instances: \textbf{in a system built out of saturating evidence, a modulation
applied upstream of the saturation is invisible.} It is worth stating as a rule, because both
occurrences were found by noticing an effect that should have been there and was not, and neither was
found by reading the code.

\subsection{No floor, and where the unresolvable go}

Any constant added underneath $P(\text{detect})$ re-creates removal-without-evidence: it guarantees
that an object never probed will eventually decay, and it does so silently, on the agent's clock. The
contract forbids it.

That leaves a real question --- an instance that is never resolvable is never reaped --- and the
answer is again the epistemic half of the formulation rather than a threshold. An unresolvable
instance is routed to a \emph{verification pull}: it raises the expected information gain of going to
look at it, the planner schedules the view, and the existence channel then decides on evidence
acquired for the purpose. Uncertainty that cannot be resolved from here becomes a reason to move,
not a reason to forget. A decay constant would have been the same decision taken without the
evidence.

Two smaller clauses complete the stage. Debouncing counts \emph{evidence cycles} rather than
wall-clock, so that mixed sensor rates do not bias removal towards whichever sensor is fastest. And
an agent may substitute the tracker's negative-information death only if it has no existence channel
at all --- a deviation that must be declared, because it changes what its deaths mean
(\S\ref{sec:attribution}).

\section{\History: no constant may stand in for experience}
\label{sec:history}

The stage that makes a model \emph{settle}, and the one where the contract is furthest ahead of both
the model and the code: nothing in \S\ref{sec:volatility} is implemented, by design, and the
recording described in \S\ref{sec:pfa} is what ships first.

The question it asks is one the belief theory of Chapter~7 does not address at all: after hours of
touring an apartment, what has the robot actually learned, and what makes a settled model resist a
spurious mask?

\subsection{The root cause is a constant}
\label{sec:constantQ}

A7 makes the state Markov, and the deployed transition for static furniture is $F=I$ with a fixed
per-frame process noise: $\Sigma\leftarrow\Sigma+Q$, applied every frame regardless of whether
anything was observed. Two consequences follow, and neither depends on how long the object has been
tracked.

\paragraph{(a) Unobserved, precision decays linearly in frames.} $\Sigma(n)=\Sigma_0+Qn$, so the
variance-doubling time is $\tau=\Sigma_0/Q$. With the shipped $\texttt{AI2ProcessStdM}=0.005$, i.e.
$Q=2.5\times10^{-5}\,\mathrm{m^2/frame}$:

\begin{center}
\begin{tabular}{@{}lrr@{}}
\toprule
& $\Sigma$ & $\tau=\Sigma/Q$ \\
\midrule
converged ($\sigma_w\approx0.03$\,m) & $9.0\times10^{-4}\,\mathrm{m^2}$ & 36 frames $\approx$ \textbf{3.6\,s} \\
prior ($\sigma_w=0.30$\,m)           & $9.0\times10^{-2}\,\mathrm{m^2}$ & 3600 frames $\approx$ 6\,min \\
\bottomrule
\end{tabular}
\end{center}

\noindent
A converged belief returns all the way to its prior after roughly $5.9$ minutes of not being looked
at, and that ceiling is the same whether the object was tracked for thirty seconds or three hours,
because nothing shrinks $Q$.

Measured, not asserted. In \texttt{table\_concept}'s log, instance \texttt{table\_9} ran cycles
$970\rightarrow2290$ --- $1320$ consecutive predict-only cycles with no assigned points throughout.
Predicted $\sigma_w=\sqrt{9.0\times10^{-4}+2.5\times10^{-5}\cdot1320}=0.184$\,m; observed
$\mathbf{0.170}$\,m, with $\sigma_{\mathrm{yaw}}$ going from $0.101$ to $0.348$\,rad over the same
stretch. The next sparse mask --- $178$ points at $2.39$\,m --- then landed on that five-times
loosened belief and re-cut the table.

\paragraph{(b) Observed, precision does not accumulate; it sits at a floor.} In information form the
update is $\Sigma^{-1}\leftarrow\Sigma^{-1}+i$, with $i$ capped at $\Lambda_c$ by
\eqref{eq:lc-sat}. The fixed point of one predict followed by one update,
$1/y=1/(y+i)+q$, is $qy^2+qiy-i=0$, giving for $qi\ll1$
\begin{equation}
y\;\approx\;\sqrt{i/q}
\qquad\Longrightarrow\qquad
\sigma_{\text{steady}}\;\approx\;\left(q/i\right)^{1/4}.
\label{eq:lc-floor}
\end{equation}
A ratio of two rates: \textbf{elapsed time does not appear}. Evidence is discarded by $Q$ at exactly
the rate it is acquired, so three hours of observation leave the belief in the same state as thirty
seconds. As an order-of-magnitude check, with $i$ at the cap $\Lambda_c=1/0.02^2=2500\,\mathrm{m^{-2}}$
and $q=2.5\times10^{-5}$, \eqref{eq:lc-floor} gives $\sigma_{\text{steady}}\approx0.01$\,m; the
observed converged $\sigma_w\approx0.03$\,m is looser because the real per-frame information sits
well below the cap once range, coverage and gating apply.

\paragraph{What is not being claimed.} $\tau=\Sigma/Q$ is \emph{not} identical across instances --- a
converged belief has the \emph{shorter} $\tau$, since a small variance doubles sooner. What is
identical is $Q$ itself, the absolute forgetting rate, and hence the retention ceiling
$\Sigma_{\text{prior}}/Q$. The load-bearing statement is that resistance to disturbance is a function
of \emph{current} $\Sigma$; current $\Sigma$ is pinned to a floor by (b) while the object is
observed, and decays to the prior by (a) when it is not, at a rate independent of history. Nothing an
agent learns extends how long it stays resistant.

This is also why clamping $\Sigma$ is not a substitute. A clamp caps the worst case of (a) and does
nothing whatever about (b). Hysteresis must be earned by evidence rather than imposed by a bound,
and recoverability must then come from something other than the bound as well.

\subsection{The move, and the first constant: volatility}
\label{sec:volatility}

Both halves of this stage are the same operation applied to two different constants:

\begin{quote}
A constant becomes an inferred, covariate-dependent latent, regularised by a hyperprior whose
pull-back keeps it recoverable.
\end{quote}

For the process noise, each instance carries a per-degree-of-freedom log-volatility $\omega$ with
$Q=\exp(\omega)$, updated after each cycle by
\begin{equation}
\omega \;\leftarrow\; \omega \;+\; \eta\left(\tfrac{1}{2}e^{-\omega}d \;-\; \tfrac{k}{2} \;-\; \frac{\omega-\omega_0}{\sigma_\omega^2}\right),
\label{eq:lc-omega}
\end{equation}
where $d=\delta\theta^{\!\top}\Sigma^{-1}\delta\theta$ is the normalised surprise of this cycle's
belief motion --- already available from the update step --- $k$ is the group's degree-of-freedom
count, and $(\omega_0,\sigma_\omega^2)$ is the hyperprior.

The test any candidate fix has to pass is that it repairs \emph{both} consequences of
\S\ref{sec:constantQ}, and \eqref{eq:lc-omega} does: $\tau\approx\Sigma e^{-\omega}$ grows as $\omega$
falls, so the retention ceiling is no longer pinned at six minutes; and
$\sigma_{\text{steady}}\approx(e^{\omega}/i)^{1/4}$ falls as $\omega$ falls, so precision genuinely
accumulates instead of resting at a fixed floor. A $\Sigma$-clamp moves only the first; an
evidence-count stiffener moves only the second. Making $Q$ itself the inferred quantity moves both,
because both are functions of $q$.

The behaviour this produces is the one the system has been approximating with tuned machinery:
consistent observation over hours keeps $d$ small, drives $\omega$ down, grows $\tau$ towards hours,
and a spurious mask can no longer move the belief --- hysteresis, earned by evidence rather than
imposed. If the object genuinely moves, sustained large $d$ raises $\omega$, $\tau$ collapses to
seconds, and the model adapts: stiffness is never a lock. Recoverability is structural rather than
thresholded, since the $(\omega-\omega_0)/\sigma_\omega^2$ term pulls $\omega$ back whenever the
evidence stops supporting extreme stiffness. And because $\omega$ is per degree of freedom, a table's
extent may freeze while its pose stays adaptive --- the correct physics for furniture, currently
asserted by hand as a zero process noise on the extent.

There is one failure mode, and it is worth flagging because it is not obvious: a persistently
\emph{biased} observation looks perfectly consistent. It drives $d$ down, drives $\omega$ down, and
stiffens the model around the bias. Volatility and the common-mode $\Sigma_\delta$ must therefore be
reasoned about together --- $\Sigma_\delta$ is exactly what prevents a shared per-frame error from
being read as agreement --- and this is the same coupling, from the other side, that
\S\ref{sec:whygate} used to justify the admissibility gate.

\subsection{The second constant: the clutter rate}
\label{sec:pfa}

The mixture in the accuracy term of \eqref{eq:lc-F} explains each point by the surface, by a wall, or
by a uniform clutter component --- and the clutter component's weight is a configured constant. That
is the second constant standing in for experience, and the same move applies.

The motivating observation is that classifiers make \emph{confident} mistakes under partial
occlusion. Part of a radiator reads as a chair with a high score, births a phantom, and the phantom
survives until a good run of masks from a proper pose kills it. Score cannot filter this, because the
score is high. But the \emph{recurrence} is learnable: the same place, seen from the same bearing,
lies the same way on every tour.

The contract therefore specifies a per-class false-alarm field $p_{\mathrm{FA}}$, indexed on
\textbf{(world cell $\times$ view bearing)}, which lowers \texttt{Detection::birth\_evidence} at
\Create. Three properties of that specification are load-bearing:

\begin{itemize}[leftmargin=1.4em,itemsep=3pt]
\item \textbf{Indexed on the pair, never on place alone.} The failure is viewpoint-dependent. A
      place-only index would suppress a \emph{genuine} object placed there from every direction;
      with bearing in the key, an unpoisoned viewpoint still births it promptly. The memory is
      never a permanent blacklist.
\item \textbf{It scales evidence and never vetoes.} The detection stays birthable; a poisoned view
      merely raises the bar, so a close, centred fixation still births a real object --- it just has
      to work harder. This is what keeps the memory falsifiable. It also reuses a hook that already
      ships, since the tracker's birth fusion moves the same scalar in the \emph{raise} direction
      from the residual field's corroboration.
\item \textbf{Same hyperprior pull-back as $\omega$}, so the field decays when unsupported. The
      world changes; a radiator can be replaced by a real chair.
\end{itemize}

\subsection{Attribution, and why recording ships first}
\label{sec:attribution}

Every agent must \emph{record}, from day one: every birth and every death, with world cell, robot
pose, view bearing, age at death, and the existence state at death --- $p_{\text{detect}}$, in field
of view, centred, fixated. Consuming the field is optional; recording is not.

The reason for that ordering is an attribution problem, and it is the same circularity
\S\ref{sec:circular} forbids in \Update. A birth followed by a death is evidence of \emph{either} a
detector false positive \emph{or} a removal false positive. Attributing it to the detector is sound
only when removal is trustworthy --- and three independent removal defects found in July 2026
(the line-of-sight threshold that deleted objects the sensor could not see, silhouettes voting
through walls per \S\ref{sec:walls}, and the inert $P(\text{detect})$ weight of
\S\ref{sec:saturation}) all produce exactly the birth-then-death signature this field would learn
from. An agent trained on those events learns its own removal bugs and calls them phantoms.

The rule that follows is stated as a hard clause: count a death as a phantom event only when the
disconfirmation was \emph{confident} --- high $p_{\text{detect}}$, fixated, close. And the error
would be self-sealing if the rule were dropped, because suppressed births generate no contradicting
evidence, so the prior becomes unfalsifiable: a wrong belief protecting itself, exactly as the
chair's obliquity covariate did.

One consequence for the audit: an agent that retires instances on \emph{divergence} rather than on
sensor absence --- because it has no existence channel at all --- has no $p_{\text{detect}}$ to
report, and its deaths are unattributable. Those events must be excluded when testing whether phantom
deaths cluster in (place $\times$ bearing), since they cannot distinguish a classifier phantom from a
model mismatch. They are marked as such in the log rather than silently mixed in.

\section{\Ownership: one mechanism, one policy}
\label{sec:ownership}

A4 partitions the object classes and gives one agent to each. That is a conditional-independence
assumption in Chapter~7 and an administrative fact here: an agent owns the graph nodes of its class,
and no other process may create or delete them.

It is worth saying why this stage is in a contract about beliefs at all. The shared graph is the only
place where the factorised posterior is reassembled --- each agent minimises its own $\sum_i F_i$ in
its own process, and what the rest of the system reads is the union of their outputs. A node whose
owner has exited is therefore a term in that product with no $q$ behind it: a factor nothing is
minimising, and one that no longer answers to any evidence. Every other stage of the contract is
about keeping the belief honest; this one is about not leaving fragments of it lying around after the
process that was maintaining it is gone.

The mechanism is a declaration in the agent's configuration --- the nodes it owns --- plus a shared
presence coordinator that removes them on shutdown. No agent may hand-roll node cleanup. It has three
ways to fail silently, and all three have been observed:

\begin{itemize}[leftmargin=1.4em,itemsep=3pt]
\item \textbf{A name that matches nothing.} Ownership entries match exact names, or a prefix when the
      entry ends in a star. An agent's own node is named after its class \emph{and its configured
      identifier}, so an entry copied from another agent --- \texttt{door\_concept 20} against a
      configured id of $15$ --- matches nothing at all, and the agent's own node leaks into the
      persistent graph on every \emph{graceful} shutdown. Found by audit, not by symptom, because a
      graceful shutdown is exactly the case nobody watches.
\item \textbf{A missing glob.} The declaration must list both the agent node and the object pattern.
      Omitting the pattern leaks every object node the agent created, or forces a bespoke sweep that
      duplicates the shared one and then drifts from it.
\item \textbf{An uncatchable signal.} Cleanup runs on \texttt{SIGTERM}/\texttt{SIGINT}. A
      \texttt{SIGKILL} cannot be caught, so nothing runs and everything owned leaks. The safety net
      is a startup stale-sweep --- run twice, because nodes may still be syncing from the persistent
      server at initialisation, so the sweep is repeated once on first entry to the operating state.
\end{itemize}

\section{What the two documents owe each other}
\label{sec:summary}

Four of the five life-cycle stages turn out to be assumptions of the factorization wearing
operational clothes, and one of them --- \History{} --- improves on the assumption as stated. That
much was the expected result of the comparison. The interesting outcome is the list of things each
document was missing that the other supplies.

\begin{center}
\begin{tabular}{@{}p{6.6cm} p{7.6cm}@{}}
\toprule
The contract gains a derivation & The model gains a correction \\
\midrule
\Update, \Remove, \Ownership{} restate the accuracy term, $q(e_i)$ and A4 --- they are not house
rules \S\ref{sec:mapping} &
The visibility factor is missing an argument: it must be conditioned on the room polygon
\eqref{eq:lc-visR} \S\ref{sec:walls} \\[5pt]
The admissibility gate is justified, not merely flagged: \eqref{eq:lc-sat} rules out the continuous
alternative \S\ref{sec:whygate} &
The gate is not a term of the likelihood but a guard on A2, whose own justification Chapter~7
records as false across short windows \S\ref{sec:whygate} \\[5pt]
Birth has a form --- the free-energy difference \eqref{eq:lc-deltaF} --- and so a currency shared
with removal \S\ref{sec:create} &
$\mathcal{I}_c$ is a fixed index set; the posterior has no birth term at all, and no birth time
\S\ref{sec:create}, \S\ref{sec:provenance} \\[5pt]
\History{} names the two constants and the single move that removes them
\S\ref{sec:volatility}, \S\ref{sec:pfa} &
A7's transition and the mixture's clutter weight are both constants that should be inferred,
covariate-dependent latents \S\ref{sec:history} \\
\bottomrule
\end{tabular}
\end{center}

\noindent
Two threads run through the list and are worth separating from it.

The first is that \eqref{eq:lc-sat} --- the one factor Chapter~7 evaluates exactly --- turned up
three times in three unrelated places: as the reason a denser mask stops helping, as the reason a
downweighted bad frame still moves the mean (\S\ref{sec:whygate}), and as the reason a carefully
computed detection weight had no effect at all (\S\ref{sec:saturation}). Saturation is the
characteristic behaviour of this architecture, and a modulation applied upstream of a saturation is
invisible. Both bugs were found by noticing a missing effect; neither was found by reading code.

The second is that the two remaining heuristics are of genuinely different kinds, which the first
draft of this chapter got wrong by grouping them. The admissibility gate is a threshold that
\emph{should} be a threshold: it defends an assumption against a bias the model has no coordinate
for, and the continuous version provably cannot do the job. The birth criterion is a threshold that
should not be: it approximates a quantity in nats with a dimensionless streak, it is the only stage
whose decision cannot be traded against any other, and closing it would subsume four separate
corrective mechanisms. Under the landmark variant of Chapter~7 it is also the stage whose errors stop
being local (\S\ref{sec:variant}). It is the first thing to fix, and the model already computes
everything it needs.

\clearpage
\appendix
\renewcommand{\thesection}{\Alph{section}}
\setcounter{section}{0}

\begin{center}
{\Large\bfseries Appendices}\\[4pt]
\rule{0.32\textwidth}{0.4pt}
\end{center}

\medskip\noindent
The body above is the contract and its reading against the model. What follows is its state on the
running system: which agents conform, where they do not, and the checklist a new agent is written
against.

\section{Conformance audit}
\label{sec:audit}

Derived by reading each agent's source, not inferred from its includes. A dash means the stage does
not apply to that agent's sensing. State as of 2026-07-31.

\begin{center}
\small
\begin{tabular}{@{}l c c c c c c c@{}}
\toprule
Agent & \Create & \Update & admissibility & \Remove & $P(\text{detect})$ & wall LoS & \Ownership \\
\midrule
table        & \ding{51} & \ding{51} & \ding{51} & \ding{51} & \ding{51} & \ding{51} & \ding{51} \\
chair        & \ding{51} & \ding{51} & \ding{51} & \ding{55}$^{a}$ & \ding{51} & --- & \ding{51} \\
bottle       & \ding{51} & \ding{51} & \ding{51} & \ding{55}$^{b}$ & --- & --- & \ding{51} \\
cabinet      & \ding{51} & \ding{51} & \ding{51} & \ding{51} & \ding{51} & \ding{51} & \ding{51} \\
door         & \ding{51} & \ding{51} & \ding{51} & \ding{51} & \ding{51} & \ding{51} & \ding{51} \\
refrigerator & \ding{51} & \ding{51} & \ding{51} & \ding{51} & \ding{51} & \ding{51} & \ding{51} \\
residual     & \ding{51} & \ding{51} & ---$^{c}$ & \ding{51} & --- & --- & \ding{51} \\
human        & \ding{55} & \ding{55} & \ding{55} & \ding{55} & \ding{55} & \ding{55} & \ding{51} \\
\bottomrule
\end{tabular}
\\[4pt]
\footnotesize
$^{a}$ hand-rolled, private accumulator \quad
$^{b}$ tracker negative-information death only \quad
$^{c}$ field model, no per-frame gate
\end{center}

\noindent
The \History{} column is omitted because it would be uninformative: \S\ref{sec:volatility} is
unimplemented in \emph{all} agents by design --- it ships after the baseline run --- and the
recording of \S\ref{sec:attribution} is live everywhere except \textsf{human\_concept}, with the
bottle and residual agents' events marked unattributable for the reason given there.

\medskip\noindent
Open deviations, in priority order:

\begin{enumerate}[leftmargin=1.6em,itemsep=3pt]
\item \textbf{\textsf{human\_concept} is on a different lineage entirely} --- no shared tracker, no
      shared belief engine, no existence channel, and the sole remaining user of the abandoned
      stabiliser module. It must either be migrated or declared out of contract explicitly; today it
      is neither, which is the worst of the three states.
\item \textbf{\textsf{chair\_concept} hand-rolls \Remove}, with a private existence accumulator and
      private gains. This is where the line-of-sight threshold of \S\ref{sec:attribution} lived. Port
      it to the shared channel.
\item \textbf{\textsf{bottle\_concept} has no existence channel}, relying on the tracker's
      negative-information death. This is permitted by \S\ref{sec:remove} but must be a conscious
      declaration, and it is why its phantom deaths are unattributable.
\item Two agents still write a legacy event schema alongside the new one. Harmless, and redundant
      once the baseline confirms the new schema captures what is wanted.
\end{enumerate}

\section{Checklist for a new agent}
\label{sec:checklist}

\begin{enumerate}[leftmargin=1.6em,itemsep=3pt]
\item \Create{} through the shared tracker; nothing else may birth. Stamp every inserted node with
      its creation time before insertion.
\item \Update{} through the shared recursive-Laplace engine. Per-frame quality into the common-mode
      covariance, never into an \texttt{if}; one stated admissibility gate, justified as in
      \S\ref{sec:whygate}; no observability covariate derived from the point estimate of the degree
      of freedom it gates.
\item \Remove{} through the shared existence channel: occupancy confirms, predicted-visible absence
      removes, occlusion holds. Absence weighted by $P(\text{detect})$ applied to the
      \emph{saturated} increment; line of sight through \texttt{walls\_block}; no floor under
      $P(\text{detect})$ --- route the unresolvable to a verification pull.
\item \History: no constant process noise once \S\ref{sec:volatility} ships; record births and
      deaths with their attribution fields from day one.
\item \Ownership: declare the agent node with the \emph{correct} identifier and the object pattern;
      terminate on catchable signals only; sweep stale nodes at startup, twice.
\item Add a row to the audit of \S\ref{sec:audit} and justify every failing entry.
\end{enumerate}

\begin{thebibliography}{9}

\bibitem{barshalom}
Y.~Bar-Shalom, X.-R.~Li and T.~Kirubarajan.
\emph{Estimation with Applications to Tracking and Navigation.} Wiley, 2001.

\bibitem{mahler}
R.~P.~S.~Mahler.
\emph{Statistical Multisource-Multitarget Information Fusion.} Artech House, 2007.

\bibitem{minka}
T.~P.~Minka.
\emph{Expectation Propagation for Approximate Bayesian Inference.}
Proc. Uncertainty in Artificial Intelligence, 2001.

\bibitem{friston}
K.~Friston, F.~Rigoli, D.~Ognibene, C.~Mathys, T.~FitzGerald and G.~Pezzulo.
\emph{Active inference and epistemic value.} Cognitive Neuroscience, 6(4):187--214, 2015.

\end{thebibliography}

\end{document}
